% ============================================================
\chapter{M\'ethodes MCMC --- Metropolis-Hastings}
\label{ch:metropolis}
% ============================================================

En 1953, au laboratoire de Los Alamos, Nicholas Metropolis, Arianna et Marshall Rosenbluth, et Augusta et Edward Teller publient un article sur la simulation de mol\'ecules par une m\'ethode de Monte Carlo particuli\`ere~: au lieu d'\'echantillonner directement la distribution de Boltzmann (impossible en haute dimension), ils construisent une cha\^ine de Markov dont la loi stationnaire \emph{est} la distribution cible. Vingt ans plus tard, W.\,Keith Hastings g\'en\'eralise l'algorithme en autorisant des noyaux de proposition asym\'etriques. Il faudra encore attendre les ann\'ees 1990 pour que les statisticiens bay\'esiens --- notamment Alan Gelfand et Adrian Smith --- r\'ealisent que cette m\'ethode r\'esout \emph{leur} probl\`eme fondamental~: \'echantillonner des lois a posteriori de forme complexe. L'algorithme de Metropolis--Hastings est ainsi devenu la cl\'e de vo\^ute de l'inf\'erence bay\'esienne computationnelle, rendant calculables des mod\`eles qui semblaient hors de port\'ee.

\begin{intuition}[title=\textbf{Id\'ee centrale}]
Quand le posterior n'a pas de forme analytique, on construit une
cha\^ine de Markov dont la loi stationnaire est la loi a posteriori.
L'algorithme de Metropolis--Hastings est la m\'ethode MCMC fondamentale~:
il propose des transitions et les accepte ou rejette selon un ratio
garantissant la convergence vers la cible.
\end{intuition}

% ------------------------------------------------------------
\section{Cha\^ines de Markov~: rappels}
% ------------------------------------------------------------

\begin{definition}[Cha\^ine de Markov]
Une suite $(\theta^{(t)})_{t \geq 0}$ est une \textbf{cha\^ine de Markov}
si $\PP(\theta^{(t+1)} \in A \mid \theta^{(0)}, \ldots, \theta^{(t)})
= \PP(\theta^{(t+1)} \in A \mid \theta^{(t)})$ pour tout $A$ mesurable.
\end{definition}

\begin{definition}[R\'eversibilit\'e et \'equilibre d\'etaill\'e]
Un noyau de transition $K(\theta, \theta')$ satisfait l'\textbf{\'equilibre
d\'etaill\'e} par rapport \`a $\pi$ si~:
\[
  \pi(\theta)\,K(\theta, \theta') = \pi(\theta')\,K(\theta', \theta)
  \quad \forall \theta, \theta'.
\]
Cela implique que $\pi$ est la loi stationnaire de la cha\^ine.
\end{definition}

\begin{theoreme}[Th\'eor\`eme ergodique]
Si la cha\^ine est irr\'eductible et ap\'eriodique avec loi
stationnaire $\pi$, alors pour toute fonction $g$ int\'egrable~:
\[
  \frac{1}{T}\sum_{t=1}^T g(\theta^{(t)}) \xrightarrow{p.s.}
  \int g(\theta)\,\pi(\theta)\,d\theta
  \quad \text{quand } T \to \infty.
\]
\end{theoreme}

% ------------------------------------------------------------
\section{Algorithme de Metropolis--Hastings}
% ------------------------------------------------------------

\begin{algorithme}[title=\textbf{Algorithme de Metropolis--Hastings}]
\begin{enumerate}
  \item Initialiser $\theta^{(0)}$.
  \item Pour $t = 0, 1, 2, \ldots, T-1$~:
  \begin{enumerate}
    \item Proposer $\theta^* \sim q(\theta^* \mid \theta^{(t)})$.
    \item Calculer le ratio d'acceptation~:
    \[
      \alpha(\theta^{(t)}, \theta^*) = \min\!\left(1,\;
      \frac{\pi(\theta^* \mid x)\,q(\theta^{(t)} \mid \theta^*)}
           {\pi(\theta^{(t)} \mid x)\,q(\theta^* \mid \theta^{(t)})}\right).
    \]
    \item Tirer $u \sim \text{Unif}(0, 1)$.
    \item Si $u \leq \alpha$~: $\theta^{(t+1)} = \theta^*$ (accepter).
          Sinon~: $\theta^{(t+1)} = \theta^{(t)}$ (rejeter).
  \end{enumerate}
\end{enumerate}
\end{algorithme}

\begin{theoreme}[Validit\'e de Metropolis--Hastings]
Le noyau de Metropolis--Hastings satisfait l'\'equilibre d\'etaill\'e
par rapport \`a $\pi(\theta \mid x)$. Si la cha\^ine est
irr\'eductible et ap\'eriodique, elle converge vers $\pi(\theta \mid x)$.
\end{theoreme}

\begin{proof}
Le noyau de transition est~:
\[
  K(\theta, \theta') = q(\theta' \mid \theta)\,\alpha(\theta, \theta')
  \quad (\theta' \neq \theta).
\]
V\'erifions l'\'equilibre d\'etaill\'e. Sans perte de g\'en\'eralit\'e,
supposons $\pi(\theta \mid x)\,q(\theta' \mid \theta) \leq
\pi(\theta' \mid x)\,q(\theta \mid \theta')$. Alors $\alpha(\theta, \theta') = 1$ et~:
\begin{align*}
  \pi(\theta \mid x)\,K(\theta, \theta')
  &= \pi(\theta \mid x)\,q(\theta' \mid \theta) \\
  &= \pi(\theta' \mid x)\,q(\theta \mid \theta')
  \cdot \frac{\pi(\theta \mid x)\,q(\theta' \mid \theta)}
             {\pi(\theta' \mid x)\,q(\theta \mid \theta')} \\
  &= \pi(\theta' \mid x)\,q(\theta \mid \theta')\,\alpha(\theta', \theta) \\
  &= \pi(\theta' \mid x)\,K(\theta', \theta).
\end{align*}
\end{proof}

\begin{remarque}
On n'a besoin que du rapport $\pi(\theta^*)/\pi(\theta^{(t)})$, donc
la constante de normalisation $m(x)$ s'annule. Il suffit de conna\^itre
le posterior \`a une constante pr\`es.
\end{remarque}

% ------------------------------------------------------------
\section{Cas particuliers}
% ------------------------------------------------------------

\subsection{Algorithme de Metropolis (sym\'etrique)}

\begin{definition}[Metropolis]
Quand la proposition est sym\'etrique, $q(\theta^* \mid \theta) = q(\theta \mid \theta^*)$,
le ratio se simplifie~:
\[
  \alpha = \min\!\left(1,\; \frac{\pi(\theta^* \mid x)}{\pi(\theta^{(t)} \mid x)}\right).
\]
\end{definition}

\begin{exemple}
\textbf{Marche al\'eatoire gaussienne.}
$q(\theta^* \mid \theta^{(t)}) = \mathcal{N}(\theta^{(t)}, \sigma_q^2 I)$.
La sym\'etrie est imm\'ediate. Le param\`etre $\sigma_q$ (pas de la
marche) doit \^etre calibr\'e~: un taux d'acceptation de $\approx 23\%$
est optimal en grande dimension.
\end{exemple}

\subsection{Algorithme d'ind\'ependance}

\begin{definition}[Metropolis--Hastings ind\'ependant]
La proposition ne d\'epend pas de l'\'etat courant~:
$q(\theta^* \mid \theta^{(t)}) = q(\theta^*)$. Le ratio est~:
\[
  \alpha = \min\!\left(1,\;
  \frac{\pi(\theta^* \mid x)\,q(\theta^{(t)})}
       {\pi(\theta^{(t)} \mid x)\,q(\theta^*)}\right).
\]
\end{definition}

\begin{attention}[title=\textbf{Condition de domination}]
L'algorithme d'ind\'ependance fonctionne bien seulement si $q$ domine
$\pi(\cdot \mid x)$ au sens o\`u les queues de $q$ sont au moins aussi
lourdes que celles de $\pi$.
\end{attention}

% ------------------------------------------------------------
\section{Calibration et adaptation}
% ------------------------------------------------------------

\begin{formules}[title=\textbf{Taux d'acceptation optimal}]
\begin{center}
\renewcommand{\arraystretch}{1.3}
\begin{tabular}{ll}
\toprule
\textbf{Dimension} & \textbf{Taux optimal} \\
\midrule
$d = 1$ & $\approx 44\%$ \\
$d \geq 2$ (cible gaussienne) & $\approx 23.4\%$ \\
\bottomrule
\end{tabular}
\end{center}
R\`egle pratique~: ajuster $\sigma_q$ pour obtenir un taux d'acceptation
entre $20\%$ et $50\%$.
\end{formules}

\begin{definition}[Metropolis adaptatif (AM)]
L'algorithme AM ajuste la matrice de covariance de la proposition
pendant le \textit{burn-in}~:
\[
  \Sigma_q^{(t)} = \frac{2.38^2}{d}\,\hat{\Sigma}_t + \varepsilon I_d,
\]
o\`u $\hat{\Sigma}_t$ est la covariance empirique des \'echantillons
$\theta^{(1)}, \ldots, \theta^{(t)}$.
\end{definition}

% ------------------------------------------------------------
\section{Diagnostic de convergence}
% ------------------------------------------------------------

\begin{definition}[Trace plot]
Le \textbf{trace plot} repr\'esente $\theta^{(t)}$ en fonction de $t$.
On recherche un comportement de \og bruit blanc \fg (bon m\'elange)
sans tendance ni corr\'elation visible.
\end{definition}

\begin{definition}[$\hat{R}$ de Gelman--Rubin]
Le diagnostic $\hat{R}$ compare la variance intra-cha\^ine $W$ et la
variance inter-cha\^ines $B$ pour $M$ cha\^ines~:
\[
  \hat{R} = \sqrt{\frac{\hat{V}}{W}}, \qquad
  \hat{V} = \frac{n-1}{n}\,W + \frac{1}{n}\,B.
\]
Convergence si $\hat{R} < 1.01$ (crit\`ere strict).
\end{definition}

\begin{definition}[Taille d'\'echantillon effective (ESS)]
L'\textbf{ESS} mesure le nombre d'\'echantillons ind\'ependants
\'equivalents~:
\[
  \text{ESS} = \frac{T}{1 + 2\sum_{k=1}^\infty \rho_k},
\]
o\`u $\rho_k$ est l'autocorr\'elation d'ordre $k$. On recommande
$\text{ESS} > 400$ par param\`etre.
\end{definition}

\begin{formules}[title=\textbf{Check-list des diagnostics MCMC}]
\begin{enumerate}
  \item V\'erifier les trace plots (pas de tendance).
  \item $\hat{R} < 1.01$ pour tous les param\`etres.
  \item $\text{ESS} > 400$ (ou $> 100$ par cha\^ine).
  \item Pas de divergences (pour HMC/NUTS).
  \item Coh\'erence entre cha\^ines multiples.
\end{enumerate}
\end{formules}

% ------------------------------------------------------------
\section{Impl\'ementation Python}
% ------------------------------------------------------------

\begin{codeblock}[title=\textbf{Metropolis--Hastings \`a la main}]
\begin{minted}{python}
import numpy as np
import matplotlib.pyplot as plt
from scipy import stats

# Cible: posterior Beta-Binomiale
n_data, s = 50, 15  # 15 succès sur 50
a_prior, b_prior = 1, 1

def log_posterior(theta):
    if theta <= 0 or theta >= 1:
        return -np.inf
    return (a_prior + s - 1) * np.log(theta) \
         + (b_prior + n_data - s - 1) * np.log(1 - theta)

# MH avec marche aléatoire
T = 50000
sigma_q = 0.1
samples = np.zeros(T)
samples[0] = 0.5
accepted = 0

for t in range(1, T):
    # Proposition
    theta_star = samples[t-1] + sigma_q * np.random.randn()
    # Ratio (Metropolis car symétrique)
    log_alpha = log_posterior(theta_star) \
              - log_posterior(samples[t-1])
    if np.log(np.random.rand()) < log_alpha:
        samples[t] = theta_star
        accepted += 1
    else:
        samples[t] = samples[t-1]

burn_in = 5000
samples_post = samples[burn_in:]
print(f"Taux d'acceptation: {accepted/T:.2%}")
print(f"Moyenne post.: {samples_post.mean():.4f}")
print(f"Analytique: {(a_prior+s)/(a_prior+b_prior+n_data):.4f}")
\end{minted}
\end{codeblock}

\begin{codeblock}[title=\textbf{Diagnostics MCMC}]
\begin{minted}{python}
fig, axes = plt.subplots(2, 2, figsize=(12, 8))

# Trace plot
axes[0, 0].plot(samples[:2000], alpha=0.7)
axes[0, 0].set_title("Trace plot (2000 premières itérations)")
axes[0, 0].set_xlabel("Itération")
axes[0, 0].set_ylabel(r"$\theta$")

# Histogramme vs analytique
theta_grid = np.linspace(0, 1, 200)
a_post, b_post = a_prior + s, b_prior + n_data - s
axes[0, 1].hist(samples_post, bins=50, density=True,
                alpha=0.5, label="MCMC")
axes[0, 1].plot(theta_grid,
                stats.beta.pdf(theta_grid, a_post, b_post),
                label="Analytique", color="red")
axes[0, 1].legend()
axes[0, 1].set_title("Posterior")

# Autocorrélation
from statsmodels.graphics.tsaplots import plot_acf
plot_acf(samples_post, lags=50, ax=axes[1, 0])
axes[1, 0].set_title("Autocorrélation")

# Running mean
running_mean = np.cumsum(samples_post) / \
               np.arange(1, len(samples_post) + 1)
axes[1, 1].plot(running_mean)
axes[1, 1].axhline((a_prior+s)/(a_prior+b_prior+n_data),
                    color="red", linestyle="--")
axes[1, 1].set_title("Moyenne cumulée")
axes[1, 1].set_xlabel("Itération")

plt.tight_layout()
plt.savefig("mcmc_diagnostics.pdf")
\end{minted}
\end{codeblock}

\begin{codeblock}[title=\textbf{MH multidimensionnel avec PyMC}]
\begin{minted}{python}
import pymc as pm
import arviz as az

# Modèle de régression bayésienne
np.random.seed(42)
n = 100
x = np.random.randn(n)
y = 2.5 + 1.3 * x + np.random.randn(n) * 0.8

with pm.Model() as model_reg:
    beta0 = pm.Normal("beta0", 0, 10)
    beta1 = pm.Normal("beta1", 0, 10)
    sigma = pm.HalfNormal("sigma", 5)
    mu = beta0 + beta1 * x
    y_obs = pm.Normal("y_obs", mu=mu, sigma=sigma,
                      observed=y)
    # Utiliser Metropolis explicitement
    step = pm.Metropolis()
    trace_mh = pm.sample(20000, step=step,
                         tune=5000, random_seed=42)

# Diagnostics
az.plot_trace(trace_mh, var_names=["beta0", "beta1", "sigma"])
plt.savefig("trace_regression.pdf")

print(az.summary(trace_mh,
                 var_names=["beta0", "beta1", "sigma"]))
\end{minted}
\end{codeblock}

% ------------------------------------------------------------
\section{Exercices}
% ------------------------------------------------------------

\begin{exercice}
Impl\'ementer l'algorithme de Metropolis avec marche al\'eatoire pour
\'echantillonner une distribution $t_3(0, 1)$. \'Etudier l'effet du
pas $\sigma_q$ sur le taux d'acceptation et l'ESS.
\end{exercice}

\begin{exercice}
D\'emontrer que l'algorithme de Metropolis--Hastings satisfait
l'\'equilibre d\'etaill\'e (preuve compl\`ete).
\end{exercice}

\begin{exercice}
Impl\'ementer l'algorithme de Metropolis adaptatif (AM) et comparer
avec Metropolis standard sur un posterior gaussien bidimensionnel avec
forte corr\'elation ($\rho = 0.95$).
\end{exercice}

\begin{exercice}
Pour un mod\`ele de m\'elange de deux gaussiennes, impl\'ementer MH et
observer le probl\`eme de multimodalit\'e. Proposer des solutions
(temp\'erage parall\`ele, propositions \`a grand saut).
\end{exercice}

\begin{exercice}
\textbf{(Num\'erique)} Comparer l'efficacit\'e (ESS par seconde) de
Metropolis--Hastings avec NUTS (PyMC par d\'efaut) sur un mod\`ele de
r\'egression logistique avec $p = 10$ pr\'edicteurs.
\end{exercice}
